% Abstract for The Lancet Digital Health submission
% LLM Proteomics Hallucination Study

\begin{abstract}

\textbf{Background:}
Large language models (LLMs) are increasingly used to answer complex biomedical questions, yet their tendency to generate plausible but incorrect information (hallucinations) poses significant risks in proteomics research and clinical applications. We evaluated hallucination rates across multiple LLMs on proteomics-specific queries.

\textbf{Methods:}
We compiled 500 unique proteomics queries spanning protein identification, quantitative expression, post-translational modifications, protein interactions, and clinical interpretation. Queries were stratified by complexity (simple/intermediate/complex) and protein prevalence (common/moderate/rare). Three frontier LLMs (GPT-4 Turbo, Claude 3 Sonnet, Gemini Pro 1.5) generated responses (1,500 total), which were evaluated by two independent expert annotators using a 4-level severity scale. Ground truth was established through multi-database validation (UniProt, Human Protein Atlas, PeptideAtlas, PhosphoSitePlus). Statistical analysis included chi-square tests with Bonferroni correction and multivariable logistic regression.

\textbf{Findings:}
Overall hallucination rate was 31.2\% (95\% CI: 28.7--33.8\%). Claude 3 Sonnet exhibited the lowest rate (27.8\%), GPT-4 Turbo intermediate (31.2\%), and Gemini Pro 1.5 highest (34.6\%, $p$<0.001 vs Claude). Hallucination rates increased markedly with query complexity (simple: 18.4\%, intermediate: 29.7\%, complex: 43.7\%; OR=5.1, 95\% CI: 4.1--6.4, $p$<0.001) and protein rarity (common: 14.3\%, moderate: 28.6\%, rare: 47.2\%; OR=5.4, 95\% CI: 4.5--6.5, $p$<0.001). Post-translational modification queries showed highest domain risk (41.8\%). Complex queries about rare proteins exceeded 50\% hallucination rates across all models. Temporal consistency was high (94.7\%, Cohen's kappa=0.91).

\textbf{Interpretation:}
Current frontier LLMs exhibit hallucination rates of 27.8--34.6\% for clinical proteomics queries, escalating to over 50\% for complex queries about rare proteins. Error rates are highest precisely where expert consultation is most critical. These hallucination rates are incompatible with safe clinical deployment. Rigorous validation frameworks and mandatory human oversight are essential before LLM deployment in proteomics applications. Regulatory evaluation standards for AI in specialized medical domains are urgently needed.

\textbf{Funding:}
This research received no external funding. Computational resources provided by DTU Computing Center.

\end{abstract}

\keywords{Large language models, hallucination, proteomics, mass spectrometry, artificial intelligence, clinical decision support, biomarkers, calibration, Bayesian analysis}
